\documentclass[12pt]{article}
\usepackage[utf8]{inputenc}
\usepackage[T5]{fontenc}
\usepackage[vietnamese]{babel}

\begin{document}

\section*{Bài tập cơ bản về Mômen lực}

\begin{enumerate}
    \item Một người dùng lực \( F = 30\,\text{N} \) để kéo một cánh cửa tại điểm cách bản lề \( d = 0{,}9\,\text{m} \). Tính mômen lực gây ra đối với bản lề.

    \item Vẫn với lực \( F = 30\,\text{N} \), nhưng người đó kéo tại điểm cách bản lề \( d = 0{,}5\,\text{m} \). Mômen lực thay đổi như thế nào so với bài 1?

    \item Một thanh gỗ dài \( 1{,}5\,\text{m} \) có điểm tựa cách đầu trái \( 0{,}5\,\text{m} \). Người A tác dụng lực \( 60\,\text{N} \) ở đầu trái. Tính lực cần đặt ở đầu phải để đòn bẩy cân bằng.

    \item Một đòn bẩy dài \( 1{,}2\,\text{m} \), điểm tựa ở giữa. Người dùng lực \( 80\,\text{N} \) ở một đầu để nâng vật ở đầu còn lại. Tính khối lượng vật nếu đòn bẩy cân bằng (lấy \( g = 9{,}8\,\text{m/s}^2 \)).

    \item Một lực \( F = 40\,\text{N} \) tác dụng lên tay quay tại góc \( 30^\circ \) so với phương vuông góc. Tay quay dài \( 0{,}6\,\text{m} \). Tính mômen lực có ích. \\ 
    \textit{Gợi ý:} Dùng công thức \( M = F \cdot r \cdot \sin(\theta) \)

    \item Hai lực \( F_1 = 20\,\text{N} \) và \( F_2 = 35\,\text{N} \) cùng tác dụng lên một thanh tại hai điểm cách trục quay lần lượt \( 0{,}4\,\text{m} \) và \( 0{,}7\,\text{m} \). Cả hai lực đều làm quay cùng chiều. Tính tổng mômen lực.

    \item Vẫn bài trên, nhưng hai lực quay ngược chiều nhau. Tính mômen lực tổng cộng. Đòn bẩy có quay không?

    \item Một lực \( 25\,\text{N} \) gây ra mômen \( 10\,\text{N}\cdot\text{m} \) quanh một trục. Tính khoảng cách từ trục quay đến điểm tác dụng.

    \item Một vật nặng \( 40\,\text{N} \) đặt ở đầu phải của thanh dài \( 1{,}2\,\text{m} \). Đặt điểm tựa cách đầu trái \( 0{,}3\,\text{m} \). Tính lực cần tác dụng ở đầu trái để giữ đòn bẩy cân bằng.

    \item Một đòn bẩy chịu hai lực: \\
    - \( F_1 = 50\,\text{N} \) tác dụng tại \( 0{,}4\,\text{m} \) về bên trái trục quay \\
    - \( F_2 = 70\,\text{N} \) tác dụng tại \( 0{,}3\,\text{m} \) bên phải trục quay \\
    Xác định chiều quay của đòn bẩy.
\end{enumerate}

\end{document}
